Physically based rainfall thresholds for unsaturated slopes usually come from uncoupled seepage analyses with
atmospheric pore air followed by stability analyses at successive instants. We quantify how these hypotheses bias the
rainfall that triggers failure. A coupled formulation with displacement, pore-water and pore-air pressures (Bishop's
effective stress, Mohr--Coulomb plasticity) is implemented in the open-source library NeoPZ and validated against the
Liakopoulos drainage test; the factor of safety (FS) is obtained by strength reduction with frozen pressures. Since
values below one are not physical states, models are compared by the critical intensity $I_c$ of a 24-h rainfall,
the smallest that brings the FS to one, and by the time to failure. For the slope of Rahardjo et al.\ (2007) with
$k_s=36$~mm/h, the uncoupled, passive-air model gives $I_c=48$~mm/h ($1.3k_s$); the coupled model with active air
gives 31~mm/h ($0.9k_s$) with an associated flow rule and 26~mm/h ($0.7k_s$) without dilatancy. The usual hypotheses
thus overestimate the failure-triggering rainfall by 54--86\% and delay failure by up to 10~h. Their effects are
opposite: trapped air halves $I_c$, coupling raises it by 4--63\%, and the air dominates. Once the surface ponds, the
air below the wetting front is compressed to 36--38~kPa; less water infiltrates and the matric suction increases, but
the pore-water pressure follows the air pressure and the effective stress drops, so that a low-permeability slope
deemed safe under short rainfalls reaches limit equilibrium. Dilatancy shifts $I_c$ by about 20\% and decides whether
the transient collapses. The bias persists with $\chi=S_e$ and mesh refinement and vanishes when $I_c$ is well below
$k_s$.
